\chapter{2013年湖南卷（理）}

\section{单选题}

\begin{problem}
复数 \(z = \i \cdot (1 + \i)\)（\(\i\) 为虚数单位）在复平面上对应的点位于（\quad）

\choices
    {第一象限}
    {第二象限}
    {第三象限}
    {第四象限}
\end{problem}

\begin{answer}
B
\end{answer}

\begin{solution}
\(z=\i(1+\i)=\i+\i^{2}=-1+\i\). 因而复数对应的点坐标为 \((-1,1)\)，实部为负、虚部为正，点位于第二象限，所以选 B.
\end{solution}

\begin{problem}
某学校有男、女学生各 \(500\text{ 名}\). 为了解男女学生在学习兴趣与业余爱好方面是否存在显著差异，拟从全体学生中抽取 \(100\text{ 名}\)学生进行调查，则宜采用的抽样方法是（\quad）

\choices
    {抽签法}
    {随机数法}
    {系统抽样法}
    {分层抽样法}
\end{problem}

\begin{answer}
D
\end{answer}

\begin{solution}
总体由男生和女生两个差异明显的层组成，且两层人数相等. 为使样本中男、女生的比例与总体一致，应在男生、女生两层中分别抽取 \(50\) 人，因此宜采用分层抽样法，选 D.
\end{solution}

\begin{problem}
在锐角 \(\triangle ABC\) 中，角 \(A\)，\(B\) 所对的边长分别为 \(a\)，\(b\). 若 \(2a \sin B = \sqrt{3} b\)，则角 \(A\) 等于（\quad）

\choices
    {\( \frac{\pi}{12} \)}
    {\( \frac{\pi}{6} \)}
    {\( \frac{\pi}{4} \)}
    {\(\frac{\pi}{3}\)}
\end{problem}

\begin{answer}
D
\end{answer}

\begin{solution}
由正弦定理，\(\dfrac{a}{\sin A}=\dfrac{b}{\sin B}\)，所以 \(a\sin B=b\sin A\). 代入题设得 \(2b\sin A=\sqrt{3}b\). 因为 \(b>0\)，故 \(\sin A=\dfrac{\sqrt{3}}{2}\). 又 \(\triangle ABC\) 为锐角三角形，\(0<A<\dfrac{\pi}{2}\)，所以 \(A=\dfrac{\pi}{3}\)，选 D.
\end{solution}

\begin{problem}
若变量 \(x\)，\(y\) 满足的约束条件 \(\begin{cases}
y \leqslant 2x,\\
x + y \leqslant 1,\\
y \geqslant -1,
\end{cases}\) 则 \(x + 2y\) 的最大值是（\quad）

\choices
    {\(-\frac{5}{2}\)}
    {0}
    {\( \frac{5}{3} \)}
    {\( \frac{5}{2} \)}
\end{problem}

\begin{answer}
C
\end{answer}

\begin{solution}
可行域由三条边界直线围成，其顶点分别为
\((-\frac{1}{2},-1)\)、\((\frac{1}{3},\frac{2}{3})\)、\((2,-1)\). 目标函数 \(x+2y\) 在这些顶点处的值依次为
\(-\frac{5}{2}\)、\(\frac{5}{3}\)、\(0\). 因为线性函数在有界多边形上的最大值在顶点处取得，所以最大值为 \(\frac{5}{3}\)，选 C.
\end{solution}

\begin{problem}
函数 \( f(x) = 2\ln x \) 的图象与函数 \( g(x) = x^2 - 4x + 5 \) 的图象的交点个数为（\quad）

\choices
    {3}
    {2}
    {1}
    {0}
\end{problem}

\begin{answer}
B
\end{answer}

\begin{solution}
令
\(h(x)=g(x)-f(x)=x^{2}-4x+5-2\ln x\)，其定义域为 \(x>0\). 有
\(h'(x)=2x-4-\dfrac{2}{x}=\dfrac{2(x^{2}-2x-1)}{x}\). 因而 \(h\) 在 \((0,1+\sqrt{2})\) 上递减，在 \((1+\sqrt{2},+\infty)\) 上递增.

又 \(\lim_{x\to0^{+}}h(x)=+\infty\)，\(\lim_{x\to+\infty}h(x)=+\infty\)，且
\(h(2)=1-2\ln2<0\)，其中 \(\ln2=\int_{1}^{2}\dfrac{1}{t}\,\mathrm{d}t>\frac{1}{2}\). 因为 \(2<1+\sqrt{2}\)，且 \(h\) 在 \((0,1+\sqrt{2})\) 上严格递减、在 \((1+\sqrt{2},+\infty)\) 上严格递增，所以在前一段由介值定理恰有一个零点，在后一段也恰有一个零点，故两图象有两个交点，选 B.
\end{solution}

\begin{problem}
已知 \(\bs{a}\)，\(\bs{b}\) 是单位向量，\(\bs{a} \cdot \bs{b} = 0\). 若向量 \(\bs{c}\) 满足 \(\left|\bs{c} - \bs{a} - \bs{b}\right| = 1\)，则 \(\left|\bs{c}\right|\) 的取值范围是（\quad）

\choices
    {\(\left[\sqrt{2} - 1, \sqrt{2} + 1\right]\)}
    {\(\left[\sqrt{2} - 1, \sqrt{2} + 2\right]\)}
    {\([1, \sqrt{2} + 1]\)}
    {\([1, \sqrt{2} + 2]\)}
\end{problem}

\begin{answer}
A
\end{answer}

\begin{solution}
由 \(|\bs a|=|\bs b|=1\) 且 \(\bs a\cdot\bs b=0\)，得
\(|\bs a+\bs b|=\sqrt{|\bs a|^{2}+|\bs b|^{2}+2\bs a\cdot\bs b}=\sqrt{2}\).
条件 \(|\bs c-\bs a-\bs b|=1\) 表明点 \(\bs c\) 在以 \(\bs a+\bs b\) 为圆心、以 \(1\) 为半径的圆上. 因此由三角不等式及其逆不等式，
\(\sqrt{2}-1\leqslant|\bs c|\leqslant\sqrt{2}+1\).
圆上沿着 \(\bs a+\bs b\) 方向和反方向的两点分别取得两个端点，所以取值范围为 \([\sqrt{2}-1,\sqrt{2}+1]\)，选 A.
\end{solution}

\begin{problem}
已知棱长为 \(1\) 的正方体的俯视图是一个面积为 \(1\) 的正方形，则该正方体的正视图的面积不可能等于（\quad）

\choices
    {1}
    {\(\sqrt{2}\)}
    {\(\frac{\sqrt{2} - 1}{2}\)}
    {\(\frac{\sqrt{2} + 1}{2}\)}
\end{problem}

\begin{answer}
C
\end{answer}

\begin{solution}
设正方体三条互相垂直的棱的单位方向向量为 \(\bs u\)，\(\bs v\)，\(\bs w\)，竖直方向的单位向量为 \(\bs n\)，并记
\(p=|\bs n\cdot\bs u|\)，\(q=|\bs n\cdot\bs v|\)，\(r=|\bs n\cdot\bs w|\). 正方体的俯视图面积等于三组棱在水平面上的投影长度之和，即
\(p+q+r=1\)，而方向余弦满足 \(p^{2}+q^{2}+r^{2}=1\). 因此
\(2(pq+qr+rp)=(p+q+r)^{2}-(p^{2}+q^{2}+r^{2})=0\). 由于 \(p\)，\(q\)，\(r\geqslant0\)，至多有一个不为零，故竖直方向平行于正方体的一条棱，正视图的竖直边长为 \(1\).

设正视方向的水平单位向量与另外两条棱的夹角分别为 \(\theta\) 和 \(\frac{\pi}{2}-\theta\)，则正视图的水平宽度为
\(w=|\cos\theta|+|\sin\theta|\). 由
\(1\leqslant w\leqslant\sqrt{2}\)，正视图面积 \(S=w\) 必满足 \(1\leqslant S\leqslant\sqrt{2}\). 选项 C 的值 \(\dfrac{\sqrt{2}-1}{2}<1\)，不可能取得，故选 C.
\end{solution}

\begin{problem}
在等腰直角三角形 \(ABC\) 中，\(AB = AC = 4\)，点 \(P\) 是边 \(AB\) 上异于 \(A\)，\(B\) 的一点，光线从点 \(P\) 出发，经 \(BC\)，\(CA\) 反射后又回到点 \(P\) （如图）. 若光线 \(QR\) 经过 \(\triangle ABC\) 的重心，则 \(AP\) 等于（\quad）

\choices
    {2}
    {1}
    {\( \frac{8}{3} \)}
    {\( \frac{4}{3} \)}
\end{problem}

\begin{answer}
D
\end{answer}

\begin{solution}
建立直角坐标系，令 \(A(0,0)\)、\(B(4,0)\)、\(C(0,4)\)，设 \(P(p,0)\)，其中 \(0<p<4\). 三角形重心为 \(G(\frac{4}{3},\frac{4}{3})\).

将点 \(P\) 关于反射边 \(BC:x+y=4\) 对称，得 \(P_{1}(4,4-p)\)；将点 \(P\) 关于反射边 \(CA:x=0\) 对称，得 \(P_{2}(-p,0)\). 反射光线展开后，点 \(P_{1}\)、反射点 \(Q\)、反射点 \(R\)、点 \(P_{2}\) 共线，所以重心 \(G\) 在直线 \(P_{1}P_{2}\) 上.

直线 \(P_{1}P_{2}\) 的方程为
\(y=\dfrac{4-p}{4+p}(x+p)\). 代入 \(G\) 得
\(\dfrac{4}{3}=\dfrac{4-p}{4+p}(\frac{4}{3}+p)\)，即
\((4-p)(3p+4)=4(4+p)\)，化简为 \(p(4-3p)=0\). 由于 \(P\) 异于 \(A\)，舍去 \(p=0\)，故 \(p=\frac{4}{3}\)，从而 \(AP=\frac{4}{3}\)，选 D.
\end{solution}

\section{填空题}

\begin{problem}
在平面直角坐标系 \(xOy\) 中，若直线 \(l: \begin{cases}
x = t,\\
y = t - a,
\end{cases}\)（\(t\) 为参数）过椭圆 \(C: \begin{cases}
x = 3\cos \varphi,\\
y = 2\sin \varphi.
\end{cases}\)（\(\varphi\) 为参数）的右顶点，则常数 \(a\) 的值为\fillinblank{}.
\end{problem}

\begin{answer}
3
\end{answer}

\begin{solution}
椭圆的右顶点对应 \(\varphi=0\)，坐标为 \((3,0)\). 直线的参数方程给出 \(x=t\)、\(y=t-a\)，代入右顶点得 \(t=3\) 且 \(3-a=0\)，所以 \(a=3\).
\end{solution}

\begin{problem}
已知 \(a\)，\(b\)，\(c \in \R\)，\(a + 2b + 3c = 6\)，则 \(a^2 + 4b^2 + 9c^2\) 的最小值为\fillinblank{}.
\end{problem}

\begin{answer}
12
\end{answer}

\begin{solution}
令 \(u=a\)、\(v=2b\)、\(w=3c\)，则 \(u+v+w=6\)，且所求式为 \(u^{2}+v^{2}+w^{2}\). 由柯西不等式，
\(u^{2}+v^{2}+w^{2}\geqslant\dfrac{(u+v+w)^{2}}{3}=12\).
当且仅当 \(u=v=w=2\)，即 \(a=2\)、\(b=1\)、\(c=\frac{2}{3}\) 时等号成立，所以最小值为 \(12\).
\end{solution}

\begin{problem}
如图，在半径为 \(\sqrt{7}\) 的 \(\odot O\) 中，弦 \(AB\)，\(CD\) 相交于点 \(P\)，\(PA = PB = 2\)，\(PD = 1\)，则圆心 \(O\) 到弦 \(CD\) 的距离为\fillinblank{}.
\bitmapfigure[max width=0.34\linewidth]{img/2013/hunan_science/q11_fig1.png}
\end{problem}

\begin{answer}
\(\frac{\sqrt{3}}{2}\)
\end{answer}

\begin{solution}
由相交弦定理，\(PA\cdot PB=PC\cdot PD\). 代入 \(PA=PB=2\)、\(PD=1\)，得 \(PC=4\)，所以弦 \(CD=PC+PD=5\).

设圆心到弦 \(CD\) 的距离为 \(d\)，连接圆心与弦中点. 由垂径定理，半弦长为 \(\frac{5}{2}\)，于是
\(d^{2}+\left(\frac{5}{2}\right)^{2}=(\sqrt{7})^{2}\). 因而 \(d=\sqrt{7-\frac{25}{4}}=\frac{\sqrt{3}}{2}\).
\end{solution}

\begin{problem}
若 \(\int_{0}^{T} x^{2} \mathrm{d}x = 9\)，则常数 \(T\) 的值为\fillinblank{}.
\end{problem}

\begin{answer}
3
\end{answer}

\begin{solution}
由微积分基本定理，\(\int_{0}^{T}x^{2}\,\mathrm{d}x=\left.\dfrac{x^{3}}{3}\right|_{0}^{T}=\dfrac{T^{3}}{3}\). 所以 \(T^{3}=27\). 因为 \(T\) 为实数，故 \(T=3\).
\end{solution}

\begin{problem}
执行如图所示的程序框图，如果输入 \(a = 1\)，\(b = 2\)，则输出的 \(a\) 的值为\fillinblank{}.
\texfigure[width=0.60\linewidth]{img_repaint/2013/hunan_science/q13_fig1.tex}
\end{problem}

\begin{answer}
9
\end{answer}

\begin{solution}
程序先输入 \(a=1\)、\(b=2\)，判断 \(a>8\). 若不成立，就把 \(b\) 加到 \(a\) 上. 各次判断前的 \(a\) 依次为 \(1\)、\(3\)、\(5\)、\(7\)、\(9\). 前四次均不大于 \(8\)，继续循环；当 \(a=9\) 时满足 \(a>8\)，输出 \(a=9\).
\end{solution}

\begin{problem}
设 \(F_{1}\)，\(F_{2}\) 是双曲线 \(C: \frac{x^{2}}{a^{2}} - \frac{y^{2}}{b^{2}} = 1 \) （\(a > 0\)，\(b > 0\)）的两个焦点，\(P\) 是 \(C\) 上一点，若 \(|PF_{1}| + |PF_{2}| = 6a\)，且 \(\triangle PF_{1}F_{2}\) 的最小内角为 \(30^{\circ}\)，则 \(C\) 的离心率为\fillinblank{}.
\end{problem}

\begin{answer}
\(\sqrt{3}\)
\end{answer}

\begin{solution}
设双曲线焦半距为 \(c\)，则 \(c^{2}=a^{2}+b^{2}\)，离心率为 \(e=\frac{c}{a}\). 双曲线的定义给出 \(\bigl||PF_{1}|-|PF_{2}|\bigr|=2a\). 与已知的距离和 \(6a\) 联立，可设两段距离分别为 \(4a\) 和 \(2a\).

三角形 \(PF_{1}F_{2}\) 的三边为 \(4a\)、\(2a\)、\(2c\). 由于 \(c>a\)，最小边为 \(2a\)，其对角为最小内角 \(30^{\circ}\). 由余弦定理，
\(4a^{2}=16a^{2}+4c^{2}-2\cdot4a\cdot2c\cos30^{\circ}\).
化简得 \(c^{2}-2\sqrt{3}ac+3a^{2}=0\)，即 \((c-\sqrt{3}a)^{2}=0\). 因而 \(e=\frac{c}{a}=\sqrt{3}\).
\end{solution}

\begin{problem}
设 \(S_{n}\) 为数列 \(\{a_{n}\}\) 的前 \(n\) 项和，\(S_{n} = (-1)^{n}a_{n} - \frac{1}{2^{n}}\)，\(n \in \N^{*}\)，则
\begin{enumerate}
\item \(a_{3} =\fillinblank{}\)；
\item \(S_{1} + S_{2} + \dots + S_{100} =\fillinblank{}\).
\end{enumerate}
\end{problem}

\begin{answer}
\(-\frac{1}{16}\)；\(\frac{1}{3}\left(\frac{1}{2^{100}} - 1\right)\)
\end{answer}

\begin{solution}
由 \(S_{1}=a_{1}\) 及题设关系（取 \(n=1\)），得 \(a_{1}=-a_{1}-\frac{1}{2}\)，所以 \(S_{1}=a_{1}=-\frac{1}{4}\).

当 \(n\) 为偶数时，\((-1)^{n}=1\)，且 \(a_{n}=S_{n}-S_{n-1}\)，代入题设关系得
\(S_{n}=S_{n}-S_{n-1}-\frac{1}{2^{n}}\)，从而 \(S_{n-1}=-\frac{1}{2^{n}}\). 因此对 \(m\geqslant1\)，有 \(S_{2m-1}=-\frac{1}{2^{2m}}\).

当 \(n\geqslant3\) 为奇数时，题设关系给出
\(S_{n}=-(S_{n}-S_{n-1})-\frac{1}{2^{n}}\)，即 \(2S_{n}=S_{n-1}-\frac{1}{2^{n}}\). 对 \(n=2m+1\)，由前面的偶数情形作用于 \(2m+2\) 得 \(S_{2m+1}=-\frac{1}{2^{2m+2}}\)，代回上式可得 \(S_{2m}=0\)（\(m\geqslant1\)）.

所以 \(a_{3}=S_{3}-S_{2}=-\frac{1}{16}\). 另外，\(S_{1}\)，\(S_{3}\)，\(\ldots\)，\(S_{99}\) 是唯一非零项，故
\(S_{1}+S_{2}+\cdots+S_{100}=-\sum_{m=1}^{50}\frac{1}{4^{m}}=-\frac{1}{3}(1-4^{-50})=\frac{1}{3}\left(\frac{1}{2^{100}}-1\right)\).
\end{solution}

\begin{problem}
设函数 \(f(x) = a^x + b^x - c^x\)，其中 \(c > a > 0\)，\(c > b > 0\).
\begin{enumerate}
\item 记集合 \(M = \{(a, b, c) \mid a, b, c\) 不能构成一个三角形的三条边长，且 \(a = b\}\)，则 \((a, b, c) \in M\) 所对应的 \(f(x)\) 的零点的取值集合为\fillinblank{}；
\item 若 \(a\)，\(b\)，\(c\) 是 \(\triangle ABC\) 的三条边长，则下列结论正确的是\fillinblank{}.（写出所有正确结论的序号）
\begin{circlelist}
\item \(\forall x \in (-\infty, 1)\)，\(f(x) > 0\)；
\item \(\exists x \in \R\)，使 \(a^{x}\)，\(b^{x}\)，\(c^{x}\) 不能构成一个三角形的三条边长；
\item 若 \(\triangle ABC\) 为钝角三角形，则 \(\exists x \in (1, 2)\)，使 \(f(x) = 0\).
\end{circlelist}
\end{enumerate}
\end{problem}

\begin{answer}
\(\left\{x \mid 0 < x \leqslant 1\right\}\)；\(\circled{1}\circled{2}\circled{3}\)
\end{answer}

\begin{solution}
\begin{enumerate}
\item 因为 \(a=b\)，且 \(c>a\)，三数不能构成三角形等价于 \(c\geqslant2a\). 令 \(t=\frac{c}{a}\)，则 \(t\geqslant2\)，且
\(f(x)=a^{x}(2-t^{x})\). 由于 \(a^{x}>0\)，零点满足 \(t^{x}=2\)，即 \(x=\frac{\ln2}{\ln t}\). 当 \(t=2\) 时 \(x=1\)，当 \(t\) 取遍 \([2,+\infty)\) 时该值取遍 \((0,1]\)，所以零点的取值集合为 \(\{x\mid0<x\leqslant1\}\).

\item 由 \(c>a\)，\(b>0\)，当 \(x<0\) 时 \(c^{x}<a^{x}\)，\(b^{x}\)，所以 \(f(x)>0\)；当 \(x=0\) 时 \(f(0)=1\). 当 \(0<x<1\) 时，三角形不等式给出 \(c<a+b\)，而幂函数 \(u^{x}\) 在正数上满足 \((a+b)^{x}<a^{x}+b^{x}\)，故 \(c^{x}<(a+b)^{x}<a^{x}+b^{x}\). 因此结论\(\circled{1}\)正确.

又因为 \(c>\max\{a,b\}\)，所以 \(0<\frac{a}{c}<1\)、\(0<\frac{b}{c}<1\). 当 \(x\) 充分大时，\(\left(\frac{a}{c}\right)^{x}+\left(\frac{b}{c}\right)^{x}<1\)，即 \(c^{x}>a^{x}+b^{x}\)，故 \(a^{x}\)，\(b^{x}\)，\(c^{x}\) 不能构成三角形，结论\(\circled{2}\)正确.

若 \(\triangle ABC\) 为钝角三角形，由于 \(c\) 是三边中最大者，钝角对边为 \(c\)，故 \(a^{2}+b^{2}<c^{2}\). 又 \(f(1)=a+b-c>0\)，\(f(2)=a^{2}+b^{2}-c^{2}<0\). 函数 \(f\) 连续，所以存在 \(x\in(1,2)\) 使 \(f(x)=0\)，结论\(\circled{3}\)正确.
\end{enumerate}
所以正确结论的序号为\(\circled{1}\circled{2}\circled{3}\).
\end{solution}

\section{解答题}

\begin{problem}
已知函数 \(f(x) = \sin \left(x - \frac{\pi}{6}\right) + \cos \left(x - \frac{\pi}{3}\right)\)，\(g(x) = 2\sin^2\frac{x}{2}\).
\begin{enumerate}
\item 若 \(\alpha\) 是第一象限角，且 \(f(\alpha) = \frac{3\sqrt{3}}{5}\)，求 \(g(\alpha)\) 的值；
\item 求使 \( f(x) \geqslant g(x) \) 成立的 \( x \) 的取值集合.
\end{enumerate}
\end{problem}

\begin{solution}
\begin{enumerate}
\item \(f(x)=\sin\left(x-\frac{\pi}{6}\right)+\cos\left(x-\frac{\pi}{3}\right)\)

\(=\frac{\sqrt{3}}{2}\sin x-\frac{1}{2}\cos x+\frac{1}{2}\cos x+\frac{\sqrt{3}}{2}\sin x=\sqrt{3}\sin x\).

\(g(x)=2\sin^{2}\frac{x}{2}=1-\cos x\).

由 \(f(\alpha)=\frac{3\sqrt{3}}{5}\) 得 \(\sin\alpha=\frac{3}{5}\). \(\alpha\) 是第一象限角，所以 \(\cos\alpha>0\)，从而

\(g(\alpha)=1-\cos\alpha=1-\sqrt{1-\sin^{2}\alpha}=1-\frac{4}{5}=\frac{1}{5}\).

\item \(f(x)\geqslant g(x)\) 等价于 \(\sqrt{3}\sin x\geqslant 1-\cos x\)，即 \(\sqrt{3}\sin x+\cos x\geqslant 1\). 于是

\(\sin\left(x+\frac{\pi}{6}\right)\geqslant\frac{1}{2}\).

从而 \(2k\pi+\frac{\pi}{6}\leqslant x+\frac{\pi}{6}\leqslant 2k\pi+\frac{5\pi}{6}\)，\(k\in\Z\)，即

\(2k\pi\leqslant x\leqslant 2k\pi+\frac{2\pi}{3}\)，\(k\in\Z\).
\end{enumerate}
\end{solution}

\begin{problem}
某人在如图所示的直角边长为 \(4\text{ 米}\) 的三角形地块的每个格点（指纵、横直线的交叉点以及三角形的顶点）处都种了一株相同品种的作物. 根据历年的种植经验，一株该种作物的年收获量 \(Y\)（单位：\(\mathrm{kg}\)）与它的“相近”作物株数 \(X\) 之间的关系如下表所示：

\begin{center}
\begin{tabular}{c|cccc}
\(X\) & 1 & 2 & 3 & 4 \\
\hline
\(Y\) & 51 & 48 & 45 & 42
\end{tabular}
\end{center}

这里，两株作物“相近”是指它们之间的直线距离不超过 \(1\text{ 米}\).

\begin{enumerate}
\item 从三角形地块的内部和边界上分别随机选取一株作物，求它们恰好“相近”的概率；
\item 从所种作物中随机选取一株，求它的年收获量的分布列与数学期望.
\end{enumerate}

\begingroup
\leftskip=0pt\relax
\rightskip=0pt\relax
\bitmapfigure[float=inline,max width=0.38\linewidth]{img/2013/hunan_science/q18_fig1.png}
\endgroup
\end{problem}

\begin{solution}
\begin{enumerate}
    \item 所种作物总数 \(N=1+2+3+4+5=15\)，其中三角形地块内部的作物株数为 \(3\)，边界上的作物株数为 \(12\). 从三角形地块的内部和边界上分别随机选取一株作物，不同结果共有 \(\mathrm C_3^1\mathrm C_{12}^1=36\) 种，选取的两株作物恰好“相近”的不同结果有 \(3+3+2=8\) 种.

把格点写成 \((i,j)\)，其中 \(i\)，\(j\geqslant0\)、\(i+j\leqslant4\). 三个内部格点为 \((1,1)\)、\((1,2)\)、\((2,1)\)，它们分别与边界上的 \(2\)、\(3\)、\(3\) 个格点相近，所以相近的结果数为 \(2+3+3=8\). 故所求概率为 \(\frac{8}{36}=\frac{2}{9}\).

\item 先求从所种作物中随机选取一株作物的年收获量 \(Y\) 的分布列. 因为

\(P(Y=51)=P(X=1)\)，\(P(Y=48)=P(X=2)\)，

\(P(Y=45)=P(X=3)\)，\(P(Y=42)=P(X=4)\).

设 \(n_k\) 为“相近”作物恰有 \(k\) 株的作物株数（\(k=1\)，\(2\)，\(3\)，\(4\)）. 逐点计数如下：顶点 \((0,0)\)、\((4,0)\)、\((0,4)\) 的相近数分别为 \(2\)，\(1\)，\(1\)；两条直角边上的其余六个格点各有 \(3\) 株相近作物；斜边上的其余三个格点各有 \(2\) 株相近作物；三个内部格点各有 \(4\) 株相近作物. 因而 \(n_1=2\)、\(n_2=4\)、\(n_3=6\)、\(n_4=3\).

由 \(P(X=k)=\frac{n_k}{15}\) 得

\(P(X=1)=\frac{2}{15}\)，\(P(X=2)=\frac{4}{15}\)，\(P(X=3)=\frac{6}{15}=\frac{2}{5}\)，\(P(X=4)=\frac{3}{15}=\frac{1}{5}\).

故所求分布列为

\begin{center}
\begin{tabular}{c|cccc}
\(Y\) & \(51\) & \(48\) & \(45\) & \(42\) \\
\hline
\(P\) & \(\frac{2}{15}\) & \(\frac{4}{15}\) & \(\frac{2}{5}\) & \(\frac{1}{5}\)
\end{tabular}
\end{center}

所求的数学期望为

\(E(Y)=51\times\frac{2}{15}+48\times\frac{4}{15}+45\times\frac{2}{5}+42\times\frac{1}{5}=\frac{34+64+90+42}{5}=46\).
\end{enumerate}
\end{solution}

\begin{problem}
如图，在直棱柱 \(ABCD - A_{1}B_{1}C_{1}D_{1}\) 中，\(AD \parallel BC\)，\(\angle BAD = 90^{\circ}\)，\(AC \perp BD\)，\(BC = 1\)，\(AD = AA_{1} = 3\).
\begin{enumerate}
\item 证明： \(AC \perp B_{1}D\)；
\item 求直线 \(B_{1}C_{1}\) 与平面 \(ACD_{1}\) 所成角的正弦值.
\end{enumerate}
\bitmapfigure[max width=0.38\linewidth]{img/2013/hunan_science/q19_fig1.png}
\end{problem}

\begin{solution}
\begin{enumerate}
\item 建立空间直角坐标系：取 \(A\) 为原点，令 \(AB\)、\(AD\)、\(AA_{1}\) 分别为 \(x\)、\(y\)、\(z\) 轴. 设 \(AB=t>0\)，由已知条件可写出
\(A(0,0,0)\)、\(B(t,0,0)\)、\(C(t,1,0)\)、\(D(0,3,0)\)、\(B_{1}(t,0,3)\)、\(C_{1}(t,1,3)\)、\(D_{1}(0,3,3)\).

于是 \(\overrightarrow{AC}=(t,1,0)\)、\(\overrightarrow{BD}=(-t,3,0)\). 由 \(AC\perp BD\)，有
\(\overrightarrow{AC}\cdot\overrightarrow{BD}=-t^{2}+3=0\)，所以 \(t=\sqrt{3}\). 此时
\(\overrightarrow{B_{1}D}=(-\sqrt{3},3,-3)\)，从而
\(\overrightarrow{AC}\cdot\overrightarrow{B_{1}D}=-3+3+0=0\)，故 \(AC\perp B_{1}D\).

\item 由上面的坐标，平面 \(ACD_{1}\) 内的两个不共线方向向量为
\(\overrightarrow{AC}=(\sqrt{3},1,0)\)、\(\overrightarrow{AD_{1}}=(0,3,3)\)，直线 \(B_{1}C_{1}\) 的方向向量为 \(\overrightarrow{B_{1}C_{1}}=(0,1,0)\).

设平面 \(ACD_{1}\) 的一个法向量为 \(n=(x,y,z)\). 则
\(\begin{cases}
n\cdot\overrightarrow{AC}=0,\\
n\cdot\overrightarrow{AD_{1}}=0,
\end{cases}\)
即
\(\begin{cases}
\sqrt{3}x+y=0,\\
3y+3z=0.
\end{cases}\)
取 \(x=1\)，得 \(n=(1,-\sqrt{3},\sqrt{3})\). 设直线 \(B_{1}C_{1}\) 与平面 \(ACD_{1}\) 所成角为 \(\theta\)，则
\(\sin\theta=\dfrac{|n\cdot\overrightarrow{B_{1}C_{1}}|}{|n|\,|\overrightarrow{B_{1}C_{1}}|}=\dfrac{\sqrt{3}}{\sqrt{7}}=\dfrac{\sqrt{21}}{7}\).
\end{enumerate}
\end{solution}

\begin{problem}
在平面直角坐标系 \(xOy\) 中，将从点 \(M\) 出发沿纵、横方向到达点 \(N\) 的任一路径称为 \(M\) 到 \(N\) 的一条“\(L\) 路径”. 如图所示的路径 \(MM_{1}M_{2}M_{3}N\) 与路径 \(MN_{1}N\) 都是 \(M\) 到 \(N\) 的“\(L\) 路径”. 某地有三个新建的居民区，分别位于平面 \(xOy\) 内三点 \(A(3,20)\)，\(B(-10,0)\)，\(C(14,0)\) 处. 现计划在 \(x\) 轴上方区域（包含 \(x\) 轴）内的某一点 \(P\) 处修建一个文化中心.
\begin{enumerate}
\item 写出点 \(P\) 到居民区 \(A\) 的“\(L\) 路径”长度最小值的表达式（不要求证明）；
\item 若以原点 \(O\) 为圆心，半径为 \(1\) 的圆的内部是保护区，“\(L\) 路径”不能进入保护区，请确定点 \(P\) 的位置，使其到三个居民区的“\(L\) 路径”长度之和最小.
\end{enumerate}
\bitmapfigure[max width=0.42\linewidth]{img/2013/hunan_science/q20_fig1.png}
\end{problem}

\begin{solution}
\begin{enumerate}
\item 设点 \(P\) 的坐标为 \((x,y)\). 点 \(P\) 到居民区 \(A\) 的“\(L\) 路径”长度最小值为

\(\lvert x-3\rvert+\lvert y-20\rvert\)，\(x\in\R\)，\(y\in[0,+\infty)\).

\item 记三个最短“\(L\) 路径”长度之和为 \(D(P)\). 不考虑保护区时，曼哈顿距离给出下界
\(D(P)\geqslant h(x)+2y+\lvert y-20\rvert\)，\(h(x)=\lvert x+10\rvert+\lvert x-14\rvert+\lvert x-3\rvert\).
由
\(\lvert x+10\rvert+\lvert x-14\rvert\geqslant24\)，且等号要求 \(x\in[-10,14]\)，再结合 \(\lvert x-3\rvert\geqslant0\)，得 \(h(x)\geqslant24\)，且等号当且仅当 \(x=3\).

当 \(y\geqslant1\) 时，若 \(1\leqslant y\leqslant20\)，则
\(D(P)\geqslant24+2y+20-y=44+y\geqslant45\). 若 \(y\geqslant20\)，则
\(D(P)\geqslant24+2y+y-20\geqslant64\).

再考虑 \(0\leqslant y<1\). 若 \(x\leqslant1\)，分段计算得：当 \(x<-10\) 时 \(h(x)=7-3x\geqslant37\)，当 \(-10\leqslant x\leqslant1\) 时 \(h(x)=27-x\geqslant26\). 所以 \(D(P)\geqslant26+20+y>45\). 若 \(x>1\)，从 \(P\) 到 \(B(-10,0)\) 的任意轴向路径必须从圆的右侧到达圆的左侧. 由于路径不能进入单位圆内部，路径在横向穿过圆的范围前必须绕到 \(y\geqslant1\) 或 \(y\leqslant-1\). 前一种走法相较于曼哈顿距离至少增加 \(2(1-y)\)；后一种走法至少增加 \(2\)，因为从高度 \(y\) 下行至 \(-1\) 再上行至 \(0\) 比竖直净位移 \(y\) 多出 \(2\). 又 \(2\geqslant2(1-y)\)，因此
\(D(P)\geqslant h(x)+20+y+2(1-y)\geqslant24+22-y>45\).

最后取 \(P=(3,1)\). 到 \(A\) 可沿直线 \(x=3\) 向上，长度为 \(19\)；到 \(B\) 先沿 \(y=1\) 向左再向下，长度为 \(14\)；到 \(C\) 先沿 \(y=1\) 向右再向下，长度为 \(12\). 这些路径均不进入保护区，长度和为 \(19+14+12=45\). 所以所求位置为 \(P(3,1)\)，最小值为 \(45\).
\end{enumerate}
\end{solution}

\begin{problem}
过抛物线 \(E: x^{2} = 2py\)（\(p > 0\)）的焦点 \(F\) 作斜率分别为 \(k_{1}\)，\(k_{2}\) 的两条不同的直线 \(l_{1}\)，\(l_{2}\)，且 \(k_{1} + k_{2} = 2\)，\(l_{1}\) 与 \(E\) 相交于点 \(A\)，\(B\)，\(l_{2}\) 与 \(E\) 相交于点 \(C\)，\(D\). 以 \(AB\)，\(CD\) 为直径的圆 \(M\)，圆 \(N\)（\(M\)，\(N\) 为圆心）的公共弦所在的直线记为 \(l\).
\begin{enumerate}
\item 若 \(k_{1} > 0\)，\(k_{2} > 0\)，证明：\(\overrightarrow{FM} \cdot \overrightarrow{FN} < 2p^{2}\)；
\item 若点 \(M\) 到直线 \(l\) 的距离的最小值为 \(\frac{7\sqrt{5}}{5}\)，求抛物线 \(E\) 的方程.
\end{enumerate}
\end{problem}

\begin{solution}
\begin{enumerate}
\item 由题意，抛物线 \(E\) 的焦点为 \(F\left(0,\frac{p}{2}\right)\)，直线 \(l_i\) 的方程为 \(y=k_i x+\frac{p}{2}\).

由 \(\begin{cases}
y=k_i x+\frac{p}{2},\\
x^2=2py
\end{cases}\) 得 \(x^2=2pk_i x+p^2\)，即

\(x^2-2pk_i x-p^2=0\).

设 \(A\)，\(B\) 两点的坐标分别为 \((x_1,y_1)\)，\((x_2,y_2)\)，则 \(x_1\)，\(x_2\) 是上述方程的两个实数根. 从而

\(x_1+x_2=2pk_1\)，

\(y_1+y_2=k_1(x_1+x_2)+p=2pk_1^2+p\).

所以 \(M\) 的坐标为 \(\left(pk_1,pk_1^2+\frac{p}{2}\right)\)，\(\overrightarrow{FM}=(pk_1,pk_1^2)\).

对直线 \(l_2\)，同样由 \(x^2=2pk_2x+p^2\) 得两交点横坐标之和为 \(2pk_2\)，纵坐标之和为 \(2pk_2^2+p\)，故
\(N=\left(pk_2,pk_2^2+\frac{p}{2}\right)\)，从而 \(\overrightarrow{FN}=(pk_2,pk_2^2)\). 于是

\(\overrightarrow{FM}\cdot\overrightarrow{FN}=p^2(k_1k_2+k_1^2k_2^2)\).

由题设，\(k_1+k_2=2\)，\(k_1>0\)，\(k_2>0\)，\(k_1\ne k_2\)，所以 \(0<k_1k_2<\left(\frac{k_1+k_2}{2}\right)^2=1\).

故 \(\overrightarrow{FM}\cdot\overrightarrow{FN}<p^2(1^2+1^2)=2p^2\).

\item 由抛物线的定义得 \(\lvert FA\rvert=y_1+\frac{p}{2}\)，\(\lvert FB\rvert=y_2+\frac{p}{2}\).

方程的根满足 \(x_1x_2=-p^2<0\)，所以焦点 \(F\) 在线段 \(AB\) 上. 因而
\(\lvert AB\rvert=\lvert FA\rvert+\lvert FB\rvert=y_1+y_2+p=2pk_1^2+2p\)，从而圆 \(M\) 的半径 \(r_1=pk_1^2+p\).

故圆 \(M\) 的方程为

\(\left(x-pk_1\right)^2+\left(y-pk_1^2-\frac{p}{2}\right)^2=\left(pk_1^2+p\right)^2\)，

化简得 \(x^2+y^2-2pk_1x-p(2k_1^2+1)y-\frac{3}{4}p^2=0\).

同样，圆 \(N\) 的圆心为 \(\left(pk_2,pk_2^2+\frac{p}{2}\right)\)，半径为 \(p(k_2^2+1)\)，所以其方程为
\(x^2+y^2-2pk_2x-p(2k_2^2+1)y-\frac{3}{4}p^2=0\).

于是圆 \(M\)，\(N\) 的公共弦所在直线 \(l\) 的方程为

\((k_2-k_1)x+(k_2^2-k_1^2)y=0\).

又 \(k_2-k_1\ne0\)，\(k_1+k_2=2\)，则 \(l\) 的方程为 \(x+2y=0\).

因为 \(p>0\)，所以点 \(M\) 到直线 \(l\) 的距离

\(d=\frac{\left|2pk_1^2+pk_1+p\right|}{\sqrt{5}}=\frac{p\left|2k_1^2+k_1+1\right|}{\sqrt{5}}=\frac{p\left[2\left(k_1+\frac{1}{4}\right)^2+\frac{7}{8}\right]}{\sqrt{5}}\).

故当 \(k_1=-\frac{1}{4}\) 时，\(d\) 取得最小值 \(\frac{7p}{8\sqrt{5}}\). 由题设 \(\frac{7p}{8\sqrt{5}}=\frac{7\sqrt{5}}{5}\)，得 \(p=8\).

故所求抛物线 \(E\) 的方程为 \(x^2=16y\).
\end{enumerate}
\end{solution}


\begin{problem}
已知 \(a > 0\)，函数 \(f(x) = \left| \frac{x - a}{x + 2a} \right|\).
\begin{enumerate}
\item 记 \(f(x)\) 在区间 \([0, 4]\) 上的最大值为 \(g(a)\)，求 \(g(a)\) 的表达式；
\item 是否存在 \(a\)，使函数 \(y = f(x)\) 在区间 \((0, 4)\) 内的图象上存在两点，在该两点处的切线相互垂直？若存在，求 \(a\) 的取值范围；若不存在，请说明理由.
\end{enumerate}
\end{problem}

\begin{solution}
\begin{enumerate}
\item 当 \(0\leqslant x\leqslant a\) 时，\(f(x)=\frac{a-x}{x+2a}\)；当 \(x>a\) 时，\(f(x)=\frac{x-a}{x+2a}\).

当 \(x\in(0,a)\) 时，\(f'(x)=\frac{-3a}{(x+2a)^2}<0\)，\(f(x)\) 在 \((0,a)\) 上单调递减；当 \(x\in(a,+\infty)\) 时，\(f'(x)=\frac{3a}{(x+2a)^2}>0\)，\(f(x)\) 在 \((a,+\infty)\) 上单调递增.

\begin{enumerate}
\item 若 \(a\geqslant4\)，则 \(f(x)\) 在 \((0,4)\) 上单调递减，\(g(a)=f(0)=\frac{1}{2}\).

\item 若 \(0<a<4\)，则 \(f(x)\) 在 \((0,a)\) 上单调递减，在 \((a,4)\) 上单调递增. 所以 \(g(a)=\max\{f(0),f(4)\}\). 而

\(f(0)-f(4)=\frac{1}{2}-\frac{4-a}{4+2a}=\frac{a-1}{a+2}\).

故当 \(0<a\leqslant1\) 时，\(g(a)=f(4)=\frac{4-a}{4+2a}\)；当 \(1<a<4\) 时，\(g(a)=f(0)=\frac{1}{2}\).
\end{enumerate}

综上所述，\(g(a)=\begin{cases}
\dfrac{4-a}{4+2a}, & 0<a\leqslant 1,\\
\dfrac{1}{2}, & a>1;
\end{cases}\)

\item 当 \(a\geqslant4\) 时，\(x<a\) 对所有 \(x\in(0,4)\) 成立，且 \(f'(x)=\dfrac{-3a}{(x+2a)^2}<0\). 所有切线斜率均为负数，任意两条切线都不可能互相垂直.

当 \(0<a<4\) 时，\(f(x)\) 在 \((0,a)\) 上单调递减，在 \((a,4)\) 上单调递增. 若存在 \(x_1\)，\(x_2\in(0,4)\)，\(x_1<x_2\)，使曲线 \(y=f(x)\) 在点 \((x_1,f(x_1))\)，\((x_2,f(x_2))\) 处的切线互相垂直，则 \(x_1\in(0,a)\)，\(x_2\in(a,4)\)，且 \(f'(x_1)\cdot f'(x_2)=-1\)，即

\(\frac{-3a}{(x_1+2a)^2}\cdot\frac{3a}{(x_2+2a)^2}=-1\).

亦即

\(x_1+2a=\frac{3a}{x_2+2a}\)，（*）

由 \(x_1\in(0,a)\)，\(x_2\in(a,4)\) 得 \(x_1+2a\in(2a,3a)\)，\(\frac{3a}{x_2+2a}\in\left(\frac{3a}{4+2a},1\right)\).

故 \((*)\) 成立，当且仅当集合 \(A=\{u\mid 2a<u<3a\}\) 与集合 \(B=\left\{u\mid\frac{3a}{4+2a}<u<1\right\}\) 有公共元素. 这里若取公共元素 \(u\)，再令 \(x_1=u-2a\)、\(x_2=\dfrac{3a}{u}-2a\)，即可得到满足范围要求的两个点，所以该条件也是充分的.

因为 \(\frac{3a}{4+2a}<3a\)，所以当且仅当 \(0<2a<1\)，即 \(0<a<\frac{1}{2}\) 时，\(A\cap B\ne\varnothing\).

综上所述，存在 \(a\) 使函数 \(y=f(x)\) 在区间 \((0,4)\) 内的图象上存在两点，在该两点处的切线互相垂直，且 \(a\) 的取值范围是 \(\left(0,\frac{1}{2}\right)\).
\end{enumerate}
\end{solution}
